\chapter{Introduction}
\label{ch:intro}

\section{Purpose }
As the demand for skilled professionals continues to grow, connecting university students with real-world opportunities has become increasingly important. Traditionally, students search for internships through fragmented processes, often relying on personal networks or generalized job boards, which may not adequately match their skills and aspirations with the needs of companies. This inefficiency results in missed opportunities for both students and employers.\newline
Students\&Companies (S\&C) is a platform designed to address this challenge by streamlining the internship matchmaking process. It bridges the gap between students seeking internships and companies offering them. S\&C allows students to showcase their experiences and skills through CVs while enabling companies to describe their projects and requirements in depth.\newline
Additionally, S\&C supports the selection process by facilitating interviews and feedback collection. With features for proactive searching and personalized suggestions, the platform aims to create a seamless and transparent experience for all parties involved.

\section{Scope}
As the RASD describes, two main users interact with the system: Students and Companies. \newline Students are individuals seeking internship opportunities. They can create and manage their profiles by uploading their CVs. The platform will notify students of relevant internship opportunities based on their profile information and allow them to apply directly. In addition, students can track their applications' status and receive company feedback. The system also provides a Recommendation Engine to suggest internships tailored to students' qualifications and preferences. Students can actively participate by proactively searching for internships or passively by checking suggested opportunities.\newline
Companies are the organizations offering internship opportunities. They will be able to post internship vacancies, specifying the required skills, tasks, technologies, and terms. Companies can manage and edit their internship postings and review and evaluate student applications. The system will help companies identify students who best match their internship requirements through a recommendation system. Companies will also be able to schedule interviews, manage the selection process, and check and respond to feedback provided by students.


\section{Definitions, Acronyms, Abbreviations}
\subsection{Definitions }
\begin{itemize}
    \item \textbf{Recommendation Engine}: System which manages the recommendation of students' CVs to companies and of internships to students
\end{itemize}

\subsection{Acronyms }
\begin{itemize}
    \item S\&C: Students\&Companies
\end{itemize}

\subsection{Abbreviations }
\begin{itemize}
    \item G*: goal
    \item WP*: world phenomena
    \item SP*: shared phenomena
    \item R*: functional requirement
    \item UC*: use case
\end{itemize}

\section{Revision history}
% Versionstabelle.

\begin{versionhistory}
	\vhEntry{1.0}{26.12.2024}{L. Ilardi, M. Luraghi}{Chapter 1}
    \vhEntry{2.0}{27.12.2024}{L. Ilardi, M. Luraghi}{...}
    \vhEntry{3.0}{28.12.2024}{L. Ilardi, M. Luraghi}{...}
    \vhEntry{4.0}{29.12.2024}{L. Ilardi, M. Luraghi}{...}
    \vhEntry{5.0}{30.12.2024}{L. Ilardi, M. Luraghi}{...}


\end{versionhistory}

\section{Reference Documents}
The document is based on the following materials:
\begin{itemize}
    \item The RASD document
    \item The specifications RASD and DD assignment of the Software Engineering II course a.a. 2024/2025
    \item Course slides on Webeep
\end{itemize}

\section{Document Structure}
\begin{enumerate}
    \item \textbf{Introduction}: it provides an overview of the project, focusing on the motivations and objectives behind its development. It also references the goals and requirements outlined in the RASD
    \item \textbf{Architectural Design}: it presents the architecture of the system, describing it at different levels of granularity
    \item \textbf{User Interface Design}: it showcases the wireframes that define the platform's user interface
    \item \textbf{Requirements Traceability}: it explains how the objectives defined in the RASD are fulfilled by the interaction between the requirements and the architectural components detailed in this document
    \item \textbf{Implementation, Integration, and Test Plan}: it describes the implementation, integration and testing details
     \item \textbf{Effort Spent}: it shows the time spent to realize this document organized by section
\end{enumerate}



\chapter{Architectural Design}
\label{Architectural Design}

\section{Overview: High-level components and their interaction}
For the S\&C platform, we will adopt a three-tier architecture, a widely used software application design that separates the system into three distinct layers:  
\begin{enumerate}
    \item \textbf{Presentation Layer}: This is the user interface, where users interact with the application.  
    \item \textbf{Application Layer}: This layer handles data processing and application logic. 
    \item \textbf{Data Layer}: This is responsible for storing and managing the application's data.   
\end{enumerate}

The primary advantage of this architecture lies in its modularity. Each layer can be developed, maintained, and scaled independently by separate teams, ensuring greater flexibility and efficiency. Moreover, changes or scaling in one tier do not directly impact the others, making the system more adaptable to evolving requirements.

\section{Component view}

\section{Deployment view}

\section{Runtime view}

\section{Component interfaces}
\textbf{InternshipManager}
\begin{itemize}
    \item getInternships() % of the company visiting the platform (and for student home?)
    \item getInternship()
    \item getStudentInternships() % the ones a student took part in or ongoing
    \item getInternshipApplications()
    \item getInternshipSuggestions()
    \item selectStudentForInternship()
    % new internship, update and delete?
\end{itemize}

\textbf{ApplicationManager}
\begin{itemize}
    \item submitApplication()
    \item getApplications() % of the student visiting the platform
    \item approveApplication()
    \item rejectApplication()
\end{itemize}

\textbf{InternshipAnnouncementManager}
\begin{itemize}
    \item getInternshipAnnouncements() % of the company visiting the platform (and for student home?)
    \item getInternshipAnnouncement()
    \item newInternshipAnnouncement()
    \item updateInternshipAnnouncement()
    \item deleteInternshipAnnouncement()
\end{itemize}

\textbf{InternshipProposalManager}
\begin{itemize}
    \item suggestInternshipToStudent()
\end{itemize}

\textbf{InterviewManager}
\begin{itemize}
    \item proposeDates()
    \item selectDate()
    \item askForReschedule()
    \item joinInterview()
\end{itemize}

\textbf{StudentProfileManager}
\begin{itemize}
    \item updateProfileInfo()
    \item uploadCv()
    \item getStudentProfile() % both for companies and the student logged in the platform
\end{itemize}

\textbf{CompanyProfileManager}
\begin{itemize}
    \item updateProfileInfo()
\end{itemize}

\textbf{FeedbackManager}
\begin{itemize}
    \item submitFeedback()
    \item answerFeedback()
\end{itemize}

\textbf{RecommendationEngine}
\begin{itemize}
    \item getStudentSuggestions()
    \item getCompanySuggestions()
\end{itemize}
\newpage

\section{Selected architectural styles and patterns}
The system follows a 3-tier architectural style due to its numerous benefits, with modularization into three independent layers being the most significant. These layers are:  

\begin{itemize}
    \item \textbf{Web Server (Presentation Tier):} This layer handles the user interface, typically presented as a website or web page. It can deliver static or dynamic content, usually developed using HTML, CSS, and JavaScript.
    \item \textbf{Application Server (Middle Tier):} This layer is responsible for implementing the business logic.
    \item \textbf{Database Server (Backend Tier):} This layer manages data storage and retrieval, operating on a database management system (DBMS) such as MySQL.
\end{itemize}

\section{Other design decisions}
\subsection{Availability }
\subsection{Data Storage }
\subsection{Security }

\newpage

\chapter{User Interface Design}
\label{User Interface Design}
%%%%%%%%%%%%%%%%%%%%%%%%%%%%%%%%%%%%%%%%%%%%%%%%%%%%%%%%%%%%%%%%%%%%%%%%%%%%%%%%%%%%%%%%%%%%%%%%
% SAME AS RASD (put in as last for readability)
%%%%%%%%%%%%%%%%%%%%%%%%%%%%%%%%%%%%%%%%%%%%%%%%%%%%%%%%%%%%%%%%%%%%%%%%%%%%%%%%%%%%%%%%%%%%%%%%
\newpage

\chapter{Requirements Traceability}
\label{Requirements Traceability}

\newpage

\chapter{Implementation, Integration and Test Plan}
\label{Implementation, Integration and Test Plan}

\section{Overview}

\section{Implementation plan}

\section{Components integration and testing}

\section{System testing}

\section{Additional specifications on testing}

\newpage

\chapter{Effort Spent}
\label{Effort Spent}